\documentclass[11pt]{article}
\usepackage[margin=1in]{geometry}
\usepackage{amsmath, amssymb}
\usepackage{graphicx}
\usepackage{booktabs}
\usepackage{float}

\title{Predicting Student Performance with Gradient Boosted Decision Trees}
\author{Michael Ellis \\ Semester Project -- CPSC 4300}
\date{November 27, 2025}

\begin{document}

\maketitle

\section{Introduction}

Supporting students who are at risk of poor academic outcomes is a
recurring challenge in educational settings.  In this project we study
the publicly available Students Performance Dataset from Kaggle, which
contains demographic information, indicators of parental and school
support, and prior achievement for several hundred students.  The goal
is to predict each student's final course performance category and to
understand which factors are most associated with high and low grades.

Following the project specification, we framed prediction as a
five--class classification problem.  The original numerical scores in
the dataset were mapped into ordered grade categories (A--F), and the
target variable $y$ represents these classes as integer labels
0--4.  Exploratory analysis of the class distribution revealed a
textbook class imbalance pattern: the majority of students fall into
the highest grade category, while very few appear in the lowest
categories.  Because of this imbalance, models that optimize overall
accuracy alone can achieve high scores by focusing primarily on the
majority class, while still performing poorly for the students we care
most about---those at risk of failing the course.

Our baseline model is an XGBoost gradient boosted decision tree
classifier, which is widely used for structured tabular data.  Using a
stratified train/validation/test split and modest hyperparameters, the
baseline XGBoost model achieved a macro--averaged F$_1$ score of
approximately 0.49 on the held--out test set.  However, inspection of
its confusion matrix and per--class metrics showed substantial
confusion among mid--range and low grades.

To address these limitations and satisfy the requirement to implement
an alternative analysis, we train a CatBoost classifier as a natural
escalation of the baseline model.  CatBoost is a gradient boosting
method specifically designed for heterogeneous tabular data with
categorical features.  It can consume categorical variables directly
without one--hot encoding and provides built--in mechanisms for
handling class imbalance.  We hypothesized that, given the strong
presence of categorical predictors (such as gender, parental
education, tutoring, and extracurricular indicators) and the skewed
class distribution, CatBoost with class--balanced loss would
outperform the XGBoost baseline in terms of macro F$_1$, particularly
for the minority grade categories.

\section{Materials and Methods}

\subsection{Dataset and Target Definition}

We use the Students Performance Dataset from Kaggle\footnote{Rabie El
Kharoua, ``Students Performance Dataset,'' Kaggle,
\\\texttt{https://www.kaggle.com/datasets/rabieelkharoua/students-performance-dataset}}.
The raw data contains a unique student identifier, demographic
variables (e.g., gender and ethnicity), indicators of parental
education and support, binary variables for tutoring and several types
of extracurricular activities, and continuous measures of prior
academic performance.

For this project, the final numeric performance score was converted
into an ordinal grade label with five ordered categories (A, B, C, D,
F).  We then mapped these categories to integer labels using a fixed
grade map
\begin{equation*}
		\text{A} \mapsto 0,\; \text{B} \mapsto 1,\; \text{C} \mapsto 2,\; \text{D} \mapsto 3,\; \text{F} \mapsto 4,
\end{equation*}
which is implemented in the project as the constant \texttt{GRADE\_MAP}.
The resulting variable, \texttt{GradeClass}, is used as the target $y$.

\subsection{Feature Engineering and Preprocessing}

All data loading and preprocessing are encapsulated in the
	exttt{reader.py} and \texttt{xgb\_baseline.py} modules.  After
loading the raw CSV file, we remove the student identifier and the
final GPA from the feature matrix because they either leak information
about the label or do not correspond to real inputs available at
prediction time.  Concretely, the model is trained on all columns
except \texttt{StudentID}, \texttt{GPA}, \texttt{GradeClass}, and
	\texttt{GradeLabel}.

The remaining predictors include both continuous variables (such as
prior exam scores) and categorical variables describing demographic
and support factors.  In the shared preprocessing function
	exttt{prepare\_data}, the following columns are treated as
categorical: \texttt{Gender}, \texttt{Ethnicity},
	\texttt{ParentalEducation}, \texttt{Tutoring},
	\texttt{ParentalSupport}, \texttt{Extracurricular}, \texttt{Sports},
	\texttt{Music}, and \texttt{Volunteering}.  For the XGBoost
baseline, these columns are cast to the \texttt{category} dtype and
fed directly to the model using its native support for categorical
features.  All remaining columns are coerced to numeric type.

The same preprocessed feature matrix $X$ and label vector $y$ are
shared between the baseline and alternative models, ensuring that
differences in performance are attributable to the modeling choices
rather than to differences in data cleaning.

\subsection{Train/Validation/Test Split}

To evaluate generalization performance, we partition the data into
training, validation, and test sets using stratified sampling on the
grade labels.  First, 20\% of the data is held out as a test set.
The remaining 80\% is further split into a training subset and a
validation subset, again using a stratified 80/20 split.  All splits
are created with a fixed random seed ($\texttt{RANDOM\_STATE} = 42$)
so that the XGBoost and CatBoost experiments share exactly the same
data partitions.

During model development, we also estimate the variability of each
model by performing 5--fold stratified cross--validation on the
training portion only.  In each fold the model is trained on 4/5 of
the training data and evaluated on the remaining 1/5.  We compute the
macro--averaged F$_1$ score in each fold and report the mean and
standard deviation across folds.

\subsection{Evaluation Metrics}

Because the grade distribution is highly imbalanced, a single overall
accuracy number can be misleading.  Instead, our primary metric is the
macro--averaged F$_1$ score on the held--out test set.  The macro
F$_1$ score gives equal weight to each grade category by averaging the
per--class F$_1$ scores,
\begin{equation*}
		\text{F}_1^{\text{macro}} = \frac{1}{K} \sum_{k=1}^K \text{F}_1^{(k)},
\end{equation*}
where $K=5$ is the number of classes.  This metric is sensitive to the
model's behaviour on rare grades and therefore better aligned with our
goal of identifying at--risk students.  For completeness we also
report overall accuracy and per--class precision, recall, and F$_1$
from the scikit--learn classification report, and we visualize the
confusion matrix for each model.

\subsection{Baseline XGBoost Model}

The baseline model is an \texttt{XGBClassifier} with the
	exttt{"multi:softprob"} objective for multi--class classification.
We use the histogram--based tree grower (\texttt{tree\_method =
"hist"}) and configure a modest set of hyperparameters that were
chosen to provide a strong yet computationally reasonable baseline:
maximum tree depth 6, learning rate 0.1, 500 boosting rounds,
subsample ratio 0.9, column subsample ratio 0.9, and L2 regularization
parameter $\lambda = 1.0$.  The model is trained with early stopping
on the validation set using the multi--class log--loss as the
evaluation metric.

We perform 5--fold stratified cross--validation on the training data
using macro F$_1$ as the scoring function.  The resulting average
macro F$_1$ score is approximately 0.57 with a standard deviation of
0.06.  On the held--out test set the XGBoost model attains an
accuracy of 0.66 and a macro F$_1$ score of 0.49.

\subsection{Alternative CatBoost Model}

The alternative model is a \texttt{CatBoostClassifier} trained on the
same feature matrix and label vector as the XGBoost baseline.  CatBoost
is particularly well suited to this dataset because it can natively
handle categorical predictors via target statistics and ordered
boosting, reducing the need for extensive feature engineering.  To
directly address the grade imbalance we enable automatic class
weighting via \texttt{auto\_class\_weights = "Balanced"}, which
increases the loss penalty for misclassifying minority classes.

We use a configuration that roughly matches the capacity of the
XGBoost baseline: depth 6 decision trees, learning rate 0.1, 500
iterations, L2 regularization parameter 3.0, and the
	\texttt{"MultiClass"} loss.  The evaluation metric during training is
CatBoost's implementation of macro F$_1$ (\texttt{TotalF1:average=Macro}).
The model is trained with early stopping using the same validation set
as the baseline model.  As with XGBoost, we estimate variability via
5--fold stratified cross--validation on the training split, recording
the macro F$_1$ score in each fold.

On the training data the CatBoost model achieves an average
cross--validated macro F$_1$ of approximately 0.61 with a standard
deviation of 0.03, already exceeding the XGBoost baseline.  When
evaluated on the held--out test set, CatBoost attains an accuracy of
0.68 and a macro F$_1$ of 0.53.

\section{Results}

In this section we compare the quantitative performance of the
baseline XGBoost model and the alternative CatBoost model, and we
summarize their error patterns and feature importances.

\subsection{Overall Performance}

Table~\ref{tab:overall-metrics} reports the cross--validated macro
F$_1$ on the training split and the accuracy and macro F$_1$ on the
held--out test set for both models.

\begin{table}[H]
	\centering
	\begin{tabular}{lccc}
		\toprule
		Model & CV F$_1^{\text{macro}}$ (mean $\pm$ std) & Test Accuracy & Test F$_1^{\text{macro}}$ \\
		\midrule
		XGBoost & 0.566 $\pm$ 0.057 & 0.658 & 0.487 \\
		CatBoost & 0.606 $\pm$ 0.035 & 0.681 & 0.532 \\
		\bottomrule
	\end{tabular}
	\caption{Overall performance of the XGBoost baseline and CatBoost
	alternative models.  All metrics are macro--averaged F$_1$ or
	accuracy on the held--out test set derived from the same
	stratified train/validation/test split.}
	\label{tab:overall-metrics}
\end{table}

Across all metrics CatBoost matches or outperforms the baseline
XGBoost model.  The cross--validated macro F$_1$ improves from 0.57 to
0.61, and the test macro F$_1$ increases from 0.49 to 0.53 while test
accuracy rises from 0.66 to 0.68.  These results support our
hypothesis that a class--balanced CatBoost model would perform at
least as well as XGBoost overall and would likely provide better
performance on the minority grade categories.

\subsection{Per--Class Metrics and Confusion Matrices}

To better understand where CatBoost improves upon XGBoost we examine
the per--class precision, recall, and F$_1$ scores from the
classification reports, as well as the confusion matrices for both
models (Figures~\ref{fig:xgb-confusion} and
\ref{fig:catboost-confusion}).

Relative to XGBoost, CatBoost yields higher F$_1$ scores for the
intermediate grade categories B, C, and D.  For example, the F$_1$
score for grade B increases from roughly 0.46 to 0.57, for grade C
from 0.47 to 0.54, and for grade D from 0.39 to 0.49.  The
improvements are driven primarily by higher recall: CatBoost correctly
identifies a larger fraction of B--, C--, and D--level students rather
than confusing them with the majority A class.  The F$_1$ score for
the failing grade F remains very high for both models, reflecting the
fact that there are relatively few failing students in the data.

The confusion matrices reinforce this pattern.  The XGBoost matrix
shows a strong concentration of predictions in the A column, indicating
that many mid--range students are over--predicted as A's.  In contrast,
CatBoost spreads its predictions more evenly across the B, C, and D
columns, reducing these systematic errors while still maintaining
strong performance on the top grade.  Figures for both models were
generated using the evaluation scripts in the project and are shown in
Figures~\ref{fig:xgb-confusion} and \ref{fig:catboost-confusion}.

\begin{figure}[H]
	\centering
	\includegraphics[width=0.48\textwidth]{"../Checkpoint 2/figures/xgb_confusion_mats.png"}%
	\includegraphics[width=0.48\textwidth]{"../data analysis/catboost_confusion_mats.png"}
	\caption{Confusion matrices for the XGBoost baseline (left) and
	CatBoost alternative (right).  Each matrix shows both raw counts
	and row--normalized values, highlighting how CatBoost reduces
	systematic misclassification of mid--range grades.}
	\label{fig:xgb-confusion}
\end{figure}

\begin{figure}[H]
	\centering
	\includegraphics[width=0.48\textwidth]{"../Checkpoint 2/figures/xgb_feature_importance.png"}%
	\includegraphics[width=0.48\textwidth]{"../data analysis/catboost_feature_importance.png"}
	\caption{Top feature importances for the XGBoost baseline (left)
	and CatBoost alternative (right).  Both models highlight prior
	performance, tutoring, and parental support as key predictors,
	with CatBoost assigning somewhat greater importance to
	categorical variables.}
	\label{fig:catboost-confusion}
\end{figure}

\subsection{Feature Importance}

We also compared the feature importance rankings produced by XGBoost
and CatBoost.  For XGBoost we used the gain--based feature importance
from the underlying booster; for CatBoost we used the model's
built--in feature importance scores.  In both cases the most
influential predictors include prior performance measures, indicators
of tutoring and parental support, and certain extracurricular
activities.  CatBoost assigns slightly higher importance to the
categorical variables, consistent with its specialized handling of
these features.  The overall similarity between the two models'
rankings provides a useful sanity check and helps interpret the
learned decision rules.

\section{Discussion and Conclusions}

The alternative CatBoost classifier successfully meets the project
requirement to implement an analysis that is anticipated to
outperform the baseline XGBoost model.  By leveraging native support
for categorical features and automatic class weighting, CatBoost
achieves higher macro--averaged F$_1$ scores and slightly better
overall accuracy on the student performance dataset.  Importantly,
the largest gains occur for the intermediate B, C, and D grade
categories, which correspond to students who are neither clearly
failing nor clearly excelling and are therefore more difficult to
classify.

From an applied perspective, these results suggest that boosting
methods that explicitly account for class imbalance and categorical
structure can provide more reliable predictions for at--risk students
than a generic gradient boosting baseline.  Nonetheless, the absolute
macro F$_1$ scores remain moderate, reflecting the inherent noise and
limited size of the dataset.  Additional improvements could likely be
obtained by exploring more extensive hyperparameter tuning, ensembling
multiple models, or incorporating richer behavioral data over time.

Despite these limitations, the project demonstrates a complete
workflow for deploying gradient boosted decision trees to predict
student performance: exploratory data analysis, careful preprocessing,
baseline modeling with XGBoost, and a targeted alternative analysis
with CatBoost that better aligns with the data's characteristics.  The
accompanying code and figures document each step of this process and
can be adapted to similar educational datasets.

\end{document}

